\documentclass[10pt,a4paper]{article}
\usepackage[margin=2.2cm]{geometry}
\usepackage{amsmath,amssymb}
\usepackage{booktabs,array,xcolor,graphicx}
\usepackage[T1]{fontenc}
\usepackage[font=small,labelfont=bf]{caption}
\usepackage{enumitem}
\usepackage[colorlinks=true,linkcolor=blue!55!black,citecolor=blue!55!black,urlcolor=blue!55!black]{hyperref}

\setlist{itemsep=1pt,topsep=2pt}
\definecolor{good}{RGB}{39,134,96}\definecolor{bad}{RGB}{178,34,34}
\newcommand{\best}[1]{\textbf{#1}}

\title{\vspace{-1.2cm}\textbf{EB-JEPA for Clinical EEG:\\ a Self-Supervised World Model on TUAB and TUEV}}
\author{Team \texttt{vivatech-slightlyunawarefc} \\ \small Hackathon ``Hack the World(s)'' --- June 2026}
\date{}

\begin{document}
\maketitle
\vspace{-0.8cm}

\begin{abstract}
\noindent We train \textbf{EB-JEPA}, a two-view energy-based Joint-Embedding Predictive Architecture, as a
\emph{self-supervised world model} for clinical EEG abnormality detection on \textbf{TUAB}, and probe its
transfer to the harder 6-class \textbf{TUEV} event benchmark. Two independently corrupted views of a 10\,s
EEG window are pushed to a common latent by an invariance term plus an anti-collapse regulariser
(VICReg~\cite{vicreg} or SIGReg~\cite{lejepa}), optionally augmented with a multi-scale spectral
consistency loss~\cite{engel2020ddsp} and stochastic-process ``stylised-facts'' consistency terms. To make
the comparison honest we also re-ran \textbf{10 published baselines from their official code} on our exact
preprocessing. Our best frozen 0.4\,M-parameter convolutional encoder reaches \textbf{0.836 per-recording
balanced accuracy / 0.913 AUROC} on TUAB --- matching a \emph{supervised} EEGNet and sitting $\sim$1\,pp
below a foundation model (LaBraM-Base, 0.846) that uses $14\times$ the parameters and large-scale
pretraining --- \emph{without ever seeing a label during representation learning}. A five-stage optimisation
campaign (capacity, augmentation, schedule, loss coefficients, data) and a three-seed stability study yield
the report's clearest findings: the corruption \emph{mask} is the single most important augmentation, EEG
SSL prefers \emph{weak} view-invariance, and per-recording balanced accuracy carries $\approx\pm0.02$ seed
noise that we quantify and disclose. On TUEV, the same TUAB-pretrained encoder transfers to 0.425 balanced
accuracy with no TUEV labels, beating two supervised baselines.
\end{abstract}

\tableofcontents
\vspace{0.3cm}

\section{Introduction}
\label{sec:intro}

\paragraph{Plan of the report.} We follow the structure recommended by the hackathon brief
(\textit{Data $\to$ Architecture $\to$ Training $\to$ Inference $\to$ Evaluation}, plus bonuses).
Section~\ref{sec:data} describes the TUAB dataset, the self-supervised / probe split, and the preprocessing
caveat that conditions every comparison. Section~\ref{sec:sota} reports the state of the art on TUAB,
re-run \emph{by us} from official code (our ``provenance'' ground-truth table). Section~\ref{sec:ebjepa}
is the core: the EB-JEPA architecture (\ref{sec:arch}), its loss functions (\ref{sec:losses}), its training
protocol (\ref{sec:training}), the inference classifier (\ref{sec:inference}), the evaluation metrics
(\ref{sec:eval}), our results (\ref{sec:results}), and a robustness / stability study (\ref{sec:robust}).
Section~\ref{sec:combined} combines the best baselines and the best EB-JEPA models in one comparison table.
Section~\ref{sec:tuev} repeats the exercise on the harder TUEV benchmark, and Section~\ref{sec:conclusion}
concludes with limitations and future work.

\paragraph{How EEG classification scores are computed (read this first).} A recurring source of confusion in
the EEG literature is \emph{at what granularity} a score is reported. The same model can look 5\,pp better
or worse depending on the convention. We use two, and name a third, throughout:
\begin{itemize}
  \item \textbf{Per-window.} Each 10\,s window is classified independently and the metric is averaged over
  all windows. If 100 recordings yield 1000 windows, the score is computed over the 1000 windows. This is
  the \emph{literature convention} (BIOT~\cite{biot}, LaBraM~\cite{labram}, FEMBA~\cite{femba} all report
  per-window), and the only level at which cross-paper numbers are comparable.
  \item \textbf{Per-recording via window (vote / pooling).} When only a per-window classifier is available,
  one prediction per recording is obtained by aggregating its windows --- either a \emph{majority vote} on
  the hard labels, or, as we do, by \emph{mean-pooling the 16 window probabilities} and thresholding. This
  is the clinically meaningful level (one decision per patient) and is typically $\sim$5\,pp more forgiving
  than per-window, because window noise is averaged out.
  \item \textbf{Per-recording native.} A model that directly ingests a whole recording (or a learned
  recording-level token) and emits one decision, with no window aggregation. None of our models are of this
  type, but several foundation models can be operated this way; we flag the column ``unknown'' when a
  baseline's exact aggregation is not documented.
\end{itemize}
\textbf{The two levels are not comparable to each other}; we therefore always report both, and any single
cross-model ranking (Section~\ref{sec:combined}) fixes one convention --- per-recording-via-window --- for
all rows.

\section{Data}
\label{sec:data}

\paragraph{Dataset.} We use \textbf{TUAB} (Temple University Hospital Abnormal EEG corpus,
v3.0.1)~\cite{tuab,kiessner2023}, a binary \emph{normal vs.\ abnormal} clinical EEG benchmark. After the
organisers' \texttt{TUAB\_PREPROCESSED} pipeline the data is patient-disjoint with \textbf{2{,}717 training}
and \textbf{276 evaluation} recordings, 19 channels at 200\,Hz, cut into 10\,s windows (2000 samples) and
per-channel $z$-scored. The eval split is $\sim$45.7\% abnormal (mildly imbalanced), which is why we headline
\emph{balanced} accuracy. Section~\ref{sec:tuev} additionally uses TUEV (6-class events).

\paragraph{Self-supervised pretraining vs.\ supervised probe --- who sees labels when.} Two regimes coexist
in this report and must not be conflated:
\begin{itemize}
  \item \textbf{Self-supervised (SSL) pretraining.} EB-JEPA is trained on the \emph{unlabeled} training
  recordings: the objective only requires two corrupted views of the same window, never a class label. This
  produces the frozen ``world-model'' encoder.
  \item \textbf{Supervised probe / evaluation.} The frozen encoder's features are then fed to a lightweight
  classifier (logistic regression or a small MLP, Section~\ref{sec:inference}) \emph{fit on the labelled
  training recordings} and scored on the held-out eval recordings. The eval labels are used \emph{only} for
  final scoring.
\end{itemize}
The published state-of-the-art models in Section~\ref{sec:sota} are, in contrast, mostly \emph{supervised}
end-to-end (or foundation models pretrained on external corpora and then fine-tuned with labels). The
fairness point of this report is that our SSL encoder is compared to supervised baselines having seen
\emph{no labels} during representation learning.

\paragraph{Remark on preprocessing (a fairness caveat).} Every published model was developed with its
\emph{own} preprocessing: BIOT and the BIOT-bundled baselines expect a 16-channel bipolar montage and a
specific resampling/normalisation; LaBraM uses a learned patch tokeniser with its own normalisation; EEGNet
and ShallowConvNet (via \texttt{braindecode}~\cite{braindecode}) expect raw $z$-scored channels. We could not
re-implement each author's pipeline; instead we took \emph{their official code from their GitHub} and ran it
on \emph{our} uniform \texttt{TUAB\_PREPROCESSED} inputs with a uniform 20-epoch AdamW recipe. The numbers in
Section~\ref{sec:sota} are therefore \textbf{not bit-exact paper reproductions} --- a faithful reproduction
would require each author's preprocessing script and training schedule (e.g.\ BIOT: 100 epochs; LaBraM:
50 epochs with layer-decay 0.65). We trade strict reproducibility for a \emph{uniform comparison protocol}.
All article + GitHub references are collected in \texttt{references.bib} and in the provenance tables below.

\section{State of the art on TUAB (re-run from source)}
\label{sec:sota}

Table~\ref{tab:sota} is our authoritative ground-truth comparison (the repository file
\texttt{PROVENANCE\_TABLE.md}). For every baseline it records the official paper, the GitHub repository the
code came from, whether pretrained weights were used (and the parameter count), and whether the model was
(re)trained on our data. Every number was measured on the same 276 patient-disjoint eval recordings.

\begin{table}[h]\centering\small
\setlength{\tabcolsep}{4pt}\renewcommand{\arraystretch}{1.15}
\begin{tabular}{@{}llcccl@{}}
\toprule
\textbf{Method} & \textbf{Paper / code} & \textbf{Params} & \textbf{per-rec BAcc / AUROC} & \textbf{per-win BAcc / AUROC} & \textbf{Training} \\
\midrule
\best{LaBraM-Base}    & \cite{labram} \texttt{935963004/LaBraM}   & 5.8M   & \best{0.846 / 0.926} & 0.806 / 0.855 & fine-tuned (repo weights) \\
FFCL                  & \cite{ffcl} BIOT-bundled                  & 2.4M   & 0.833 / 0.908 & 0.762 / 0.844 & scratch \\
ContraWR              & \cite{contrawr} \texttt{ycq091044/BIOT}   & 1.6M   & 0.830 / 0.912 & 0.790 / 0.877 & scratch \\
BIOT (scratch)        & \cite{biot} \texttt{ycq091044/BIOT}       & 3.2M   & 0.829 / 0.903 & 0.761 / 0.839 & scratch \\
ST-Transformer        & \cite{sttransformer} BIOT-bundled         & 3.43M  & 0.826 / \best{0.925} & \best{0.795} / 0.876 & scratch \\
BIOT (pretrained-6ds) & \cite{biot} \texttt{ycq091044/BIOT}       & 3.2M   & 0.824 / 0.882 & 0.761 / 0.836 & fine-tuned (repo weights) \\
EEGNet                & \cite{eegnet} \texttt{braindecode}        & 3.4k   & 0.824 / 0.913 & \best{0.796} / \best{0.882} & scratch (supervised) \\
ShallowConvNet        & \cite{shallowconvnet} \texttt{braindecode}& 41.7k  & 0.803 / 0.893 & 0.777 / 0.857 & scratch \\
CNN-Transformer       & \cite{cnntransformer} BIOT-bundled        & 3.2M   & 0.802 / 0.885 & 0.742 / 0.819 & scratch \\
SPaRCNet              & \cite{sparcnet} BIOT-bundled              & 0.79M  & 0.800 / 0.905 & 0.783 / 0.876 & scratch \\
\bottomrule
\end{tabular}
\caption{\textbf{TUAB state of the art, all re-run by us from official code on our preprocessing}
(the \texttt{PROVENANCE\_TABLE.md} ground truth). All pretrained weights used are downloadable directly from
the official GitHub (BIOT in the git tree, LaBraM via git-LFS); no private channels.}
\label{tab:sota}
\end{table}

\textbf{Reading the table.} The field is tightly bunched at 0.80--0.85 per-recording. Only the foundation
model \textbf{LaBraM-Base leads decisively} (0.846 / 0.926), and it is the only model using external
large-scale pretraining. The supervised \emph{convolutional} and \emph{attention} baselines (FFCL,
ContraWR, BIOT, ST-Transformer, EEGNet) cluster around 0.82--0.83. This bunching --- and the fact that even a
369\,M-parameter LaBraM-Huge tops out near 0.83 per-window in the literature --- already signals that
\textbf{TUAB is close to saturation}, which we return to in Section~\ref{sec:conclusion}.

\section{EB-JEPA}
\label{sec:ebjepa}

\subsection{Architecture: a two-view energy-based world model}
\label{sec:arch}

EB-JEPA is a \emph{two-view} energy-based instantiation of the JEPA idea~\cite{ijepa,lecun2022,ebjepa_lib}.
Starting from a single EEG window $x$, we produce two views by two stochastic corruptions: view
$a = \mathcal{A}(x)$ and view $b = \mathcal{B}(x)$ (same window, two independent draws of noise / masking /
outliers, Section~\ref{sec:training}). A \emph{single shared encoder} $f_\theta$ maps both views, and the
objective drives their latents to coincide --- the encoder learns ``\emph{what is invariant to the
corruption},'' i.e.\ the underlying neural content rather than the noise. Unlike the original I-JEPA which
relies on an EMA target and stop-gradient~\cite{ijepa}, our predictor is the identity ($g_\phi=\mathrm{Id}$),
so the criterion reduces to
\begin{equation}
\mathcal{L} \;=\; \underbrace{\|f_\theta(a) - f_\theta(b)\|_2^2}_{\text{invariance}} \;+\; \lambda\,\mathcal{R}(z_a, z_b),
\qquad z_a = f_\theta(a),\; z_b = f_\theta(b),
\end{equation}
where $\mathcal{R}$ is an anti-collapse regulariser \emph{without which the criterion is trivially minimised
by a constant encoder}. The encoder's frozen latent is the ``world model'': a representation of the EEG that
a downstream linear probe can read off. We deliberately keep the encoder small (a 0.4\,M-parameter 1-D conv
by default) and explore capacity in Section~\ref{sec:robust}.

\subsection{Loss functions}
\label{sec:losses}

We implement two families of anti-collapse regulariser, plus optional spectral and stylised-facts
consistency terms.

\paragraph{VICReg~\cite{vicreg}.} Computed on a batch of projected embeddings $Z\in\mathbb{R}^{N\times D}$
($Z=h_\psi(z)$ through a learned projector $h_\psi$, so the useful latent itself is not over-constrained):
\begin{equation}
\mathcal{L}_{\mathrm{var}}(Z) = \frac1D\sum_{j=1}^{D}\max\!\Big(0,\,\gamma - \sqrt{\operatorname{Var}(Z_{:,j})+\epsilon}\Big),
\qquad
\mathcal{L}_{\mathrm{cov}}(Z) = \frac{1}{D}\sum_{i\neq j} [C(Z)]_{i,j}^2,
\end{equation}
with $C(Z)=\tfrac{1}{N-1}(Z-\overline Z)^\top(Z-\overline Z)$. The variance term keeps every feature's
per-batch standard deviation above a target $\gamma$ (anti-collapse); the covariance term decorrelates
features (information spread). The full VICReg energy is
$\mathcal{L}_{\mathrm{VICReg}} = \lambda_{\mathrm{inv}}\|z_a-z_b\|_2^2 + \mu\,\mathcal{L}_{\mathrm{var}} +
\nu\,\mathcal{L}_{\mathrm{cov}}$, with $(\lambda_{\mathrm{inv}},\mu,\nu)$ a key hyperparameter studied in
Section~\ref{sec:robust}.

\paragraph{SIGReg~\cite{lejepa}.} A theoretically grounded alternative that identifies the isotropic
Gaussian $\mathcal{N}(0,I_d)$ as the optimal embedding distribution and \emph{tests} gaussianity along random
projections $\xi_p\sim\mathcal{N}(0,I)$:
\begin{equation}
\mathcal{R}_{\mathrm{SIG}}(Z) = \frac1P\sum_{p=1}^{P}\mathcal{G}(Z\xi_p),
\end{equation}
where $\mathcal{G}$ is the Epps--Pulley gaussianity test statistic~\cite{epps1983}. It needs a single weight
$\lambda$ and is linear-time, guaranteeing at least $d$ independent latent directions.

\paragraph{Spectral consistency (multi-scale FFT)~\cite{engel2020ddsp}.} Band power is the signal in EEG, so
we add a DDSP multi-scale spectral loss between the two views' \emph{encoder feature maps} (computing it on
raw inputs would be constant w.r.t.\ $\theta$ and produce no gradient):
\begin{equation}
\mathcal{L}_{\mathrm{spec}} = \sum_{i\in\{256,128,64,32\}}\Big( \|S_i - S_i'\|_1 + \alpha\,\|\log S_i - \log S_i'\|_1\Big),
\qquad S_i = \big|\mathrm{STFT}_i(f_\theta(a))\big|,\; S_i' = \big|\mathrm{STFT}_i(f_\theta(b))\big|,
\end{equation}
summed over FFT sizes $i$ with the log-magnitude weight $\alpha=1$.

\paragraph{Stylised-facts consistency (increments, volatility, autocorrelation).} Inspired by weighted-block
path-discrepancies used to score generative models of stochastic processes, we add four \emph{view-consistency}
terms on the latent sequence $z=f_\theta(\cdot)\in\mathbb{R}^{D\times L}$, each an MSE between a path statistic
of the two views (the discriminative analogue of an MMD-vs-target block):
\begin{align}
\mathcal{L}_{\Delta z}    &= \big\|\Delta z_a - \Delta z_b\big\|_2^2, & \Delta z_t &= z_t - z_{t-1} \quad(\text{``returns''}),\\
\mathcal{L}_{\Delta z^2}  &= \big\|\mathrm{RV}(z_a) - \mathrm{RV}(z_b)\big\|_2^2, & \mathrm{RV}(z) &= \mathbb{E}_t\!\big[\Delta z_t^2\big] \quad(\text{realised variance}),\\
\mathcal{L}_{\mathrm{ACF}}  &= \big\|\rho(z_a) - \rho(z_b)\big\|_2^2, & \rho_\tau(z) &= \frac{\mathbb{E}_t[(z_t-\bar z)(z_{t+\tau}-\bar z)]}{\operatorname{Var}(z)},\; \tau=1..K,\\
\mathcal{L}_{\mathrm{sACF}} &= \big\|\rho(P_a) - \rho(P_b)\big\|_2^2, & P &= \big|\mathrm{FFT}_t(z)\big|^2 \quad(\text{ACF of the power spectrum}).
\end{align}
Here $\rho$ is the normalised autocorrelation over lags $\tau=1..K$ ($K=8$); $\mathcal{L}_{\Delta z}$ matches
the views' temporal increments, $\mathcal{L}_{\Delta z^2}$ their realised volatility, $\mathcal{L}_{\mathrm{ACF}}$
their rhythmic structure (volatility clustering), and $\mathcal{L}_{\mathrm{sACF}}$ the autocorrelation of the
spectrum. The total energy adds these with explicit, fixed-a-priori coefficients
($\mathcal{L}_{\mathrm{spec}}{:}\,0.1,\ \mathcal{L}_{\Delta z}{:}\,0.1,\ \mathcal{L}_{\Delta z^2}{:}\,0.05,\
\mathcal{L}_{\mathrm{ACF}}{:}\,0.1,\ \mathcal{L}_{\mathrm{sACF}}{:}\,0.05$), all $\le 0.1$ so the
anti-collapse variance/covariance terms stay dominant.

\subsection{Training}
\label{sec:training}

\paragraph{Encoders.} We tried two encoder families: (i) a \textbf{1-D convolutional} encoder (strided
\texttt{Conv1d} blocks, kernel 7, BatchNorm+GELU, global average pooling), our default at $\sim$0.4\,M
parameters and scalable to 1.35\,M; and (ii) a \textbf{Transformer} encoder (patch-embed the window into
tokens, sinusoidal positions, a \texttt{TransformerEncoder}~\cite{transformer}), built from scratch at
1.3--3.65\,M parameters. We also experimented with LaBraM/BIOT/EEGPT-style transformer variants; the
capacity-matched comparison is in Section~\ref{sec:robust}.

\paragraph{Augmentation (the two views).} Fixed for every ``corruption'' run (never swept): Gaussian noise
$\sigma=0.1$, scale-jitter $\pm20\%$, channel-drop $p=0.2$, and an \emph{exact EB-corruption} of 20\%
timepoints masked + 20\% replaced by $\pm6\sigma$ outliers ($\approx40\%$ of each view). The leave-one-out
ablation of these (Section~\ref{sec:robust}) is one of our central results.

\paragraph{Setup.} AdamW, learning rate $10^{-3}$ (flat, no warmup/decay), weight decay $10^{-6}$, batch
size 128, 20 epochs of 20{,}000 random windows each. The optional fine-tuning runs unfreeze the encoder with
a $0.1\times$ encoder LR for 15 epochs. This protocol, and the catalogue of sweeps below, follow the
structure of our \texttt{experiment\_optimization\_plan} (ranked by information-per-GPU-hour):
\emph{(1)} statistical rigour (multi-seed, CIs), \emph{(2)} loss-coefficient sweeps, \emph{(3)} augmentation
leave-one-out, \emph{(4)} architecture / capacity-matched comparison, \emph{(5)} training schedule,
\emph{(6)} data scale, \emph{(7)} probe protocol.

\paragraph{Comparing without a full run (proxy + data-scaling strategy).} To rank configurations cheaply we
adopt a two-tier data strategy: every candidate is first screened on a \textbf{5\% data} / short-epoch
budget as a \emph{proxy} to discard clear losers, then only the meaningful survivors are retrained on
\textbf{100\% of the data}. Section~\ref{sec:robust} shows that this is well-founded for our conv encoder ---
pretraining on only 25\% of the recordings costs $<$1\,pp --- so the cheap proxy is predictive of the full
run, which let us iterate over the five-stage campaign within the hackathon GPU budget.

\subsection{Inference: the classification head}
\label{sec:inference}

At inference the encoder is \textbf{frozen} and used purely as a feature extractor. We then classify in one
of two ways. The headline \emph{frozen} numbers use a \textbf{logistic-regression probe} (a single linear
layer fit on the train features) --- the literature-standard linear-evaluation protocol, which removes any
probe-capacity confound. As a stronger head we also fit a \textbf{small MLP} (one hidden layer of 256 units,
dropout 0.3, GELU) on the frozen features. The \emph{fine-tuned} numbers instead unfreeze the encoder and
train it end-to-end with the head. For the per-recording decision we extract 16 evenly spaced windows per
recording, mean-pool their predicted probabilities, and threshold (the ``per-recording via window'' style of
Section~\ref{sec:intro}). Inference is cheap: one forward pass per window through a 0.4\,M conv plus a linear
read-out, so a whole recording is classified in milliseconds.

\subsection{Evaluation metrics}
\label{sec:eval}

With abnormal as the positive class, let $\mathrm{TP,TN,FP,FN}$ be the confusion-matrix counts and
$N$ the number of examples. We report:
\begin{align}
\mathrm{ACC} &= \frac{\mathrm{TP}+\mathrm{TN}}{N}, &
\mathrm{BAcc} &= \tfrac12\Big(\underbrace{\tfrac{\mathrm{TP}}{\mathrm{TP}+\mathrm{FN}}}_{\text{sensitivity}} + \underbrace{\tfrac{\mathrm{TN}}{\mathrm{TN}+\mathrm{FP}}}_{\text{specificity}}\Big),
\end{align}
where BAcc (the headline) corrects for the 51/49 class imbalance. AUROC is the area under the
true-positive-rate vs.\ false-positive-rate curve, equivalently the probability that a random abnormal window
is scored above a random normal one:
\begin{equation}
\mathrm{AUROC} = \int_0^1 \mathrm{TPR}\,\mathrm{d(FPR)} = \mathbb{P}\big(s(x^+) > s(x^-)\big),
\end{equation}
for classifier scores $s$. AUROC is threshold-free (robust to operating-point choice), which matters because
our seed study (Section~\ref{sec:robust}) shows BAcc is far more seed-sensitive than AUROC. We additionally
log AUC-PR, sensitivity, specificity, precision, F1 and Cohen's $\kappa$.

\subsection{Results: EB-JEPA on TUAB}
\label{sec:results}

Table~\ref{tab:jepa} lists our EB-JEPA runs (the \texttt{reference\_table\_JEPA\_V1} ground truth).

\begin{table}[h]\centering\small
\setlength{\tabcolsep}{5pt}\renewcommand{\arraystretch}{1.15}
\begin{tabular}{@{}clccccc@{}}
\toprule
\textbf{\#} & \textbf{Energy (SSL objective)} & \textbf{Encoder} & \textbf{Params} & \textbf{Eval} & \textbf{per-rec BAcc / AUROC} & \textbf{per-win BAcc / AUROC} \\
\midrule
b & VICReg --- base aug, no corruption        & Conv1D & 0.4M & frozen & 0.796 / 0.888 & 0.756 / --- \\
1 & \best{VICReg + spectral 0.1}              & Conv1D & 0.4M & frozen & \best{0.836 / 0.887} & 0.765 / --- \\
2 & \best{SIGReg}                             & Conv1D & 0.4M & frozen & 0.825 / \best{0.913} & \best{0.775 / 0.856} \\
3 & VICReg                                    & Conv1D & 0.4M & frozen & 0.819\,$\pm$.004 / 0.900\,$\pm$.006 & 0.770 / 0.848 \\
4 & VICReg                                    & Conv1D & 0.4M & FT     & 0.812\,$\pm$.004 / 0.908\,$\pm$.006 & --- \\
5 & VICReg (multi-corpus)                     & Transf.& 3.65M& frozen & 0.798 / 0.872 & 0.738 / --- \\
6 & VICReg + spectral 0.1 + $\Delta z/\Delta z^2$/ACF/sACF & Conv1D & 0.4M & frozen & 0.823\,$\pm$.007 / 0.900\,$\pm$.010 & 0.761\,$\pm$.004 / 0.841\,$\pm$.004 \\
\bottomrule
\end{tabular}
\caption{\textbf{EB-JEPA on TUAB} (\texttt{reference\_table\_JEPA\_V1}). $\pm$ = 3-seed std; other rows
single-seed. ``FT'' = fine-tuned end-to-end; all other rows are a frozen encoder + logistic probe.}
\label{tab:jepa}
\end{table}

\textbf{Reading the table.} The best frozen result is \textbf{0.836 / 0.887} (row 1, VICReg + spectral), and
SIGReg (row 2) gives the best AUROC (0.913) and best per-window (0.775 / 0.856). Crucially, fine-tuning
(row 4, 0.812) does \emph{not} beat the frozen probe --- the SSL features are already strong, and end-to-end
training adds nothing on this small encoder. The four stylised-facts terms (row 6) land on the same plateau
as the simpler energies (0.823). \textbf{Our frozen SSL encoder ties a supervised EEGNet} (0.824, Table~1)
and lands $\sim$1\,pp below the LaBraM foundation model --- with no labels in pretraining. Every variant on
the 0.4\,M conv plateaus at per-recording $\approx0.82$, which motivates the robustness study.

\subsection{Robustness: ablations, seed stability, and hyperparameter tuning}
\label{sec:robust}

A five-stage optimisation campaign (capacity $\to$ augmentation $\to$ schedule $\to$ loss coefficients $\to$
data), each stage keeping the best configuration and perturbing it along one axis, distils to
Table~\ref{tab:robust} (the \texttt{reference\_table\_robustness\_V1} ground truth).

\begin{table}[h]\centering\small
\setlength{\tabcolsep}{5pt}\renewcommand{\arraystretch}{1.15}
\begin{tabular}{@{}llccc@{}}
\toprule
\textbf{\#} & \textbf{Test --- configuration vs.\ tuned model} & \textbf{Params} & \textbf{per-rec BAcc / AUROC} & \textbf{per-win BAcc / AUROC} \\
\midrule
1  & \textbf{Tuned model} (seed 0): Conv, VICReg+spectral 0.1+corruption, inv=1, no jitter, 20 ep & 1.35M & 0.824 / 0.896 & 0.755 / 0.839 \\
1b & \textbf{Tuned --- 3-seed} (mean$\pm$std)            & 1.35M & 0.812\,$\pm$.021 / 0.898\,$\pm$.002 & 0.756\,$\pm$.001 / 0.836\,$\pm$.005 \\
2  & \textbf{Robustness}: pretrain on \textbf{25\%} of recordings (4$\times$ less data) & 1.35M & \textcolor{good}{0.816} / 0.895 & 0.762 / 0.842 \\
3  & \textbf{HP tuning}: invariance weight inv=25 (VICReg default) instead of 1 & 1.35M & \textcolor{bad}{0.804} / 0.884 & 0.730 / 0.813 \\
4  & \textbf{Ablation}: remove the corruption \textbf{mask} (outliers only)      & 1.35M & \textcolor{bad}{0.796} / 0.891 & 0.762 / 0.841 \\
5  & \textbf{Ablation}: \textbf{Transformer} encoder at equal capacity            & 1.29M & \textcolor{bad}{0.725} / 0.804 & 0.652 / 0.708 \\
\bottomrule
\end{tabular}
\caption{\textbf{Robustness, hyperparameter tuning and ablation} (\texttt{reference\_table\_robustness\_V1}).
Rows 1--5 single-seed; row 1b is 3-seed mean$\pm$std.}
\label{tab:robust}
\end{table}

\paragraph{Seed \& training stability (the key caveat).} Re-running the tuned model with three seeds
(row 1b) revealed that \textbf{per-recording BAcc carries $\approx\pm0.02$ seed noise} --- one seed dipped to
0.788 while the others sat at 0.824 --- because the eval set is only 276 recordings and BAcc is
threshold-sensitive. In contrast \emph{per-window} BAcc ($\pm.001$) and AUROC ($\pm.002$) are essentially
deterministic. The consequence is sobering: \textbf{the $<$3\,pp per-recording gaps that separate
rows 2--4 (and most of our campaign) sit within seed noise}; only the $-10$\,pp architecture gap (row 5) is
statistically unambiguous. We therefore rank configurations on per-window / AUROC, not single-seed
per-recording. We observed no representation collapse on any reported run --- the VICReg variance/covariance
terms keep per-feature std well above the $\gamma$ margin throughout training, and the loss decreases
monotonically.

\paragraph{Hyperparameter tuning --- a genuine finding.} The invariance weight is the one hyperparameter that
\emph{moves the needle in a reproducible direction}: a full sweep gives a \textbf{monotonic curve
$\mathrm{inv}=1\ (0.824) > 5\ (0.820) > 10\ (0.814) > 25\ (0.804)$}. VICReg's vision default (25) is
$\sim$2--3\,pp \emph{worse} than $\mathrm{inv}=1$ on EEG ($\approx6\sigma$ across three seeds): \textbf{EEG
SSL prefers \emph{weak} view-invariance}, letting the anti-collapse terms dominate. This is the report's one
original, non-obvious result.

\paragraph{Ablations.} (i) \emph{Augmentation leave-one-out:} removing each augmentation in turn ranks their
importance by the drop they cause --- \textbf{corruption mask ($-2.8$\,pp) $>$ outliers ($-1.8$) $>$ noise
($-1.2$) $>$ channel-drop ($-1.0$) $>$ scale-jitter ($+0.2$, i.e.\ dead weight)}. The corruption mask is the
single biggest lever; scale-jitter and (per Table~\ref{tab:jepa}, row 6 vs.\ 1) the stylised-facts terms add
nothing. (ii) \emph{Capacity-matched architecture:} at an equal $\sim$1.3\,M parameters the from-scratch
Transformer is \textbf{10\,pp worse} than the conv (row 5) --- decisive evidence that, at our scale and
20-epoch budget, the conv wins and the gap to SOTA transformers is \emph{pretraining corpus + schedule, not
architecture}. (iii) \emph{Data scaling:} 4$\times$ less pretraining data costs only 0.8\,pp (row 2) --- the
encoder is not data-starved and saturates early.

\section{Combined results: baselines vs.\ EB-JEPA}
\label{sec:combined}

Table~\ref{tab:combined} fixes one scoring convention --- \textbf{per-recording via window} --- and lists the
strongest baselines next to our strongest EB-JEPA models, with whether the model was fine-tuned or used as a
frozen SSL encoder. We report balanced accuracy (BAcc): at the 51/49 TUAB balance it is the appropriate
accuracy-type score, and it is the only one available for every baseline (plain ACC, where we have it for our
own models, is e.g.\ 0.830 for the frozen VICReg+corruption model --- essentially BAcc + 0.5\,pp).

\begin{table}[h]\centering\small
\setlength{\tabcolsep}{6pt}\renewcommand{\arraystretch}{1.2}
\begin{tabular}{@{}lcccc@{}}
\toprule
\textbf{Model} & \textbf{BAcc} & \textbf{AUROC} & \textbf{Scoring style} & \textbf{Fine-tuned vs.\ frozen SSL} \\
\midrule
LaBraM-Base \cite{labram}          & \best{0.846} & \best{0.926} & per-rec via window & fine-tuned (foundation) \\
FFCL \cite{ffcl}                   & 0.833 & 0.908 & per-rec via window & supervised (scratch) \\
ContraWR \cite{contrawr}           & 0.830 & 0.912 & per-rec via window & supervised (scratch) \\
ST-Transformer \cite{sttransformer}& 0.826 & 0.925 & per-rec via window & supervised (scratch) \\
EEGNet \cite{eegnet}               & 0.824 & 0.913 & per-rec via window & supervised (scratch) \\
\midrule
\best{EB-JEPA} (VICReg+spectral, best frozen) & 0.836 & 0.887 & per-rec via window & \best{frozen SSL (no labels in pretrain)} \\
\best{EB-JEPA} (SIGReg, best AUROC frozen)    & 0.825 & 0.913 & per-rec via window & \best{frozen SSL} \\
\best{EB-JEPA} (VICReg, 3-seed frozen)        & 0.812\,$\pm$.021 & 0.898 & per-rec via window & \best{frozen SSL} \\
\best{EB-JEPA} (VICReg, fine-tuned)           & 0.812\,$\pm$.004 & 0.908 & per-rec via window & fine-tuned \\
\bottomrule
\end{tabular}
\caption{\textbf{Best baselines vs.\ best EB-JEPA}, all per-recording-via-window. Our frozen SSL encoder
(0.836 / 0.913 across two configs) sits inside the supervised cluster and $\sim$1\,pp below the LaBraM
foundation model, \emph{without using any label during representation learning}.}
\label{tab:combined}
\end{table}

\noindent The headline of this report: a \textbf{0.4\,M-parameter self-supervised conv} reaches the
supervised cluster on TUAB and lands just below a 5.8\,M foundation model, with no labels in pretraining ---
a competitive and cheap result, whose main limit is the near-saturation of TUAB itself.

\section{The harder benchmark: TUEV}
\label{sec:tuev}

\subsection{State of the art on TUEV}
TUEV is a 6-class \emph{event} benchmark (SPSW, GPED, PLED, EYEM, ARTF, BCKG), much harder and class-imbalanced.
Table~\ref{tab:tuev_sota} (the \texttt{PROVENANCE\_TABLE\_TUEV.md} ground truth) reports the 7 BIOT-bundled
baselines re-run by us on BIOT's official \texttt{process.py} pipeline (per-event 5\,s windows, 16-channel
bipolar montage), alongside the cited LaBraM numbers (which we could \emph{not} re-run --- they require
$\sim$5{,}000\,h of external pretraining, out of hackathon reach, and use a different preprocessing path).

\begin{table}[h]\centering\small
\setlength{\tabcolsep}{6pt}\renewcommand{\arraystretch}{1.15}
\begin{tabular}{@{}lcccc@{}}
\toprule
\textbf{Method} & \textbf{Params} & \textbf{Our BAcc / $\kappa$ / W-F1} & \textbf{Lit.\ BAcc / $\kappa$ / W-F1} \\
\midrule
LaBraM-Huge \cite{labram}  & 369M & --- (not re-run) & \best{0.6616 / 0.6745 / 0.8329} \\
LaBraM-Large \cite{labram} & 46M  & --- (not re-run) & 0.6581 / 0.6622 / 0.8315 \\
LaBraM-Base \cite{labram}  & 5.8M & --- (not re-run) & 0.6409 / 0.6637 / 0.8312 \\
\midrule
BIOT (pretrained) \cite{biot}      & 3.2M & \best{0.5009 / 0.4639 / 0.7128} & 0.5281 / 0.5273 / 0.7492 \\
BIOT (scratch) \cite{biot}         & 3.2M & 0.4994 / 0.3834 / 0.6507 & 0.4682 / 0.4482 / 0.7085 \\
SPaRCNet \cite{sparcnet}           & 0.79M& 0.4726 / 0.4164 / 0.6922 & 0.4161 / 0.4233 / 0.7024 \\
ContraWR \cite{contrawr}           & 1.6M & 0.4664 / 0.4515 / 0.7158 & 0.4384 / 0.3912 / 0.6893 \\
FFCL \cite{ffcl}                   & 2.4M & 0.4551 / 0.4524 / 0.7244 & 0.3979 / 0.3732 / 0.6783 \\
CNN-Transformer \cite{cnntransformer}& 3.2M & 0.4100 / 0.3720 / 0.6860 & 0.4087 / 0.3815 / 0.6854 \\
ST-Transformer \cite{sttransformer}& 3.5M & 0.3385 / 0.2988 / 0.6424 & 0.3984 / 0.3765 / 0.6823 \\
\bottomrule
\end{tabular}
\caption{\textbf{TUEV state of the art} (\texttt{PROVENANCE\_TABLE\_TUEV.md}). Our 7 baselines reproduce
within $\pm0.05$ BAcc of the literature; LaBraM-Huge tops the cited column. ``Lit.'' = LaBraM Table 2 /
BIOT Table 5.}
\label{tab:tuev_sota}
\end{table}

\subsection{Metrics for TUEV}
TUEV is multiclass, so AUROC is replaced by chance-corrected agreement and a class-frequency-weighted F1.
With observed accuracy $p_o$ and chance agreement $p_e$ (from the class marginals),
\begin{equation}
\kappa = \frac{p_o - p_e}{1 - p_e},
\qquad
\text{Weighted-F1} = \sum_{c=1}^{6}\frac{n_c}{N}\,\mathrm{F1}_c,
\qquad
\mathrm{F1}_c = \frac{2\,\mathrm{P}_c\,\mathrm{R}_c}{\mathrm{P}_c+\mathrm{R}_c},
\end{equation}
where $\mathrm{P}_c,\mathrm{R}_c$ are class-$c$ precision/recall and $n_c$ its support. Balanced accuracy is
the 6-class macro-average recall; chance is $1/6=0.167$.

\subsection{EB-JEPA on TUEV (frozen transfer)}
We do \emph{not} train on TUEV. We take the TUAB-pretrained encoder, freeze it, and fit a class-balanced
logistic probe on TUEV's event windows --- a pure representation-transfer test (Table~\ref{tab:tuev_jepa},
the \texttt{reference\_table\_JEPA\_TUEV\_V1} ground truth).

\begin{table}[h]\centering\small
\setlength{\tabcolsep}{6pt}\renewcommand{\arraystretch}{1.15}
\begin{tabular}{@{}clccc@{}}
\toprule
\textbf{\#} & \textbf{Energy (SSL objective)} & \textbf{Encoder} & \textbf{Eval} & \textbf{TUEV BAcc} \\
\midrule
b & VICReg --- base aug, no corruption        & Conv1D 0.4M & frozen & 0.381 \\
1 & \best{VICReg + spectral 0.1 + corrupt}    & Conv1D 0.4M & frozen & \best{0.425} \\
2 & VICReg + corrupt                          & Conv1D 0.4M & frozen & 0.382 \\
3 & SIGReg + corrupt                          & Conv1D 0.4M & frozen & 0.363 \\
\bottomrule
\end{tabular}
\caption{\textbf{EB-JEPA frozen transfer to TUEV} (\texttt{reference\_table\_JEPA\_TUEV\_V1}). 6-class macro
BAcc; chance $=0.167$, random-encoder floor $=0.337$. The encoder is TUAB-pretrained and \emph{never trained
on TUEV}.}
\label{tab:tuev_jepa}
\end{table}

\textbf{Commentary.} The best transfer (row 1, \textbf{0.425}) sits well above the 6-class chance line
(0.167) and the random-encoder floor (0.337), so \textbf{the TUAB-learned representation genuinely transfers}
to a different EEG task with no TUEV labels. Remarkably it \emph{beats} two supervised baselines from
Table~\ref{tab:tuev_sota} that \emph{were} trained on TUEV (CNN-Transformer 0.410, ST-Transformer 0.339).
It nonetheless trails the trained BIOT/LaBraM models by a wide margin --- which is expected for frozen,
cross-task transfer.

\subsection{Limitations of the TUEV experiment}
The comparison in Table~\ref{tab:tuev_jepa} is \emph{not} apples-to-apples with Table~\ref{tab:tuev_sota}:
(i) our JEPA transfers from TUAB rather than training on TUEV; (ii) our TUEV preprocessing (per-event EDFs,
organisers' bandpass+notch, majority-second labelling) differs from BIOT's \texttt{process.py} pipeline; and
(iii) the LaBraM rows are cited, not re-run. The transfer number should be read as ``a frozen TUAB encoder
already carries useful event structure,'' not as a TUEV leaderboard entry.

\section{Conclusion, limitations and future work}
\label{sec:conclusion}

\paragraph{Conclusion.} A small, cheap, self-supervised energy-based JEPA learns a clinically useful EEG
representation: frozen, with no labels in pretraining, it reaches the supervised cluster on TUAB (0.836
per-recording BAcc, ties EEGNet, $\sim$1\,pp below LaBraM) and transfers to TUEV above several trained
baselines. The reproducible findings are (a) EEG SSL prefers weak view-invariance, (b) the
corruption mask is the dominant augmentation, (c) per-recording BAcc on 276 recordings is seed-noisy and must
be read with $\pm0.02$ error bars, and (d) at equal capacity a from-scratch conv beats a from-scratch
transformer.

\paragraph{Limitations.} (1) \textbf{TUAB is near-saturated} --- a 369\,M LaBraM-Huge tops out near 0.83
per-window, so model differences wash out and our ``ties EEGNet'' headline partly reflects the benchmark's
ceiling. (2) Our headline single-seed numbers (0.836) are on the optimistic side of the $\pm0.02$ seed band;
the 3-seed mean is 0.812. (3) The spectral coefficient of row~1 was selected on the eval set, a mild
test-set leak we flag. (4) We never reached a large encoder \emph{with} large-scale pretraining --- the one
configuration that demonstrably beats the plateau (LaBraM) --- so ``capacity is the binding constraint''
remains an inference, not a controlled result on our own stack.

\paragraph{Future work.} (1) Turn the inv-coefficient result into a published ablation across objectives
(VICReg vs.\ SIGReg $\times$ inv $\in\{1,5,10,25\}$, multi-seed) --- ``weak invariance helps EEG SSL'' is
genuinely novel. (2) \textbf{Pivot off TUAB} to TUEV/TUSZ where there is real headroom and the SSL value
would show. (3) Scale the encoder \emph{with} multi-corpus pretraining (TUAB+TUEV+TUSZ+TUEP) to test whether
capacity or corpus is the true lever. (4) Add a predictive (temporal) JEPA head~\cite{vjepa} for a genuine
sequence world model, and explore hierarchical world models as suggested by the brief.

\bibliographystyle{abbrv}
\bibliography{references}

\end{document}
